\documentclass[conference]{IEEEtran}
\IEEEoverridecommandlockouts
% The preceding line is only needed to identify funding in the first footnote. If that is unneeded, please comment it out.
\usepackage{cite}
\usepackage{amsmath,amssymb,amsfonts}
\usepackage{algorithmic}
\usepackage{graphicx}
\usepackage{textcomp}
\usepackage{xcolor}
\usepackage{fancyhdr}
\usepackage{hyperref}

\def\BibTeX{{\rm B\kern-.05em{\sc i\kern-.025em b}\kern-.08em
    T\kern-.1667em\lower.7ex\hbox{E}\kern-.125emX}}

\fancypagestyle{firstpage}{
  \fancyhf{} % Clear header and footer
  \fancyhead[C]{23rd Edition of IEEE India Council International Conference (INDICON 2026)}
  \fancyfoot[L]{%
      \begin{minipage}[t]{\textwidth}
      \normalsize{979-8-3195-2581-9/26/\$31.00 \copyright2026 IEEE}
    \end{minipage}%
  }
  \renewcommand{\headrulewidth}{0pt} % Remove the horizontal line
}

\pagestyle{fancy}
\fancyhf{} % Clear the default header and footer
\fancyfoot[C]{INDICON 2026}
\renewcommand{\headrulewidth}{0pt} % Remove the horizontal line

\begin{document}

\title{Conference Paper Title for 23rd Edition of IEEE India Council International Conference}

\author{\IEEEauthorblockN{Given Name Surname}
\IEEEauthorblockA{\textit{dept. name of organization (of Aff.)} \\
\textit{name of organization (of Aff.)}\\
City, Country \\
email address }
\and
\IEEEauthorblockN{Given Name Surname}
\IEEEauthorblockA{\textit{dept. name of organization (of Aff.)} \\
\textit{name of organization (of Aff.)}\\
City, Country \\
email address }
\and
\IEEEauthorblockN{Given Name Surname\textsuperscript{*}}
\IEEEauthorblockA{\textit{dept. name of organization (of Aff.)} \\
\textit{name of organization (of Aff.)}\\
City, Country \\
email address }
\and
\IEEEauthorblockN{Given Name Surname}
\IEEEauthorblockA{\textit{dept. name of organization (of Aff.)} \\
\textit{name of organization (of Aff.)}\\
City, Country \\
email address }
\and
\IEEEauthorblockN{Given Name Surname}
\IEEEauthorblockA{\textit{dept. name of organization (of Aff.)} \\
\textit{name of organization (of Aff.)}\\
City, Country \\
email address }
\and
\IEEEauthorblockN{Given Name Surname}
\IEEEauthorblockA{\textit{dept. name of organization (of Aff.)} \\
\textit{name of organization (of Aff.)}\\
City, Country \\
email address }
}

\maketitle
\thispagestyle{firstpage}

\begin{abstract}
The 23rd Edition of the IEEE India Council International Conference (INDICON 2026) provides an international forum for researchers, academicians, industry professionals, and practitioners to present and discuss recent advances in engineering, computing, communication technologies, and emerging interdisciplinary fields. Organized by the IEEE India Council in association with the IEEE Madras Section, the conference focuses on the theme “Net-Zero Cyber-Physical Intelligence: AI, 6G \& Sustainable Electronics for Future-Ready Smart Societies.” The conference encompasses a broad range of research areas, including artificial intelligence and machine learning, next-generation wireless communication, cyber-physical systems, digital twins, sustainable electronics, very-large-scale integration (VLSI), robotics, intelligent control, photonics, quantum technologies, signal processing, biomedical engineering, and information security. INDICON 2026 aims to foster knowledge exchange, encourage interdisciplinary collaboration, and promote innovative technological solutions to address emerging societal and sustainability challenges. The conference will be held from December 18–20, 2026, at Sri Sairam Engineering College, Chennai, India, in a hybrid format, enabling both in-person and virtual participation.
\end{abstract}

\begin{IEEEkeywords}
IEEE INDICON 2026, artificial intelligence, cyber-physical systems, 6G communications, sustainable electronics, net-zero engineering.
\end{IEEEkeywords}

\section{Introduction}
The 23rd Edition of the IEEE India Council International Conference (INDICON 2026) brings together the global research community to explore emerging technologies and their role in shaping sustainable, intelligent, and future-ready societies. Organized by the IEEE India Council in association with the IEEE Madras Section, the conference is guided by the theme “Net-Zero Cyber-Physical Intelligence: AI, 6G \& Sustainable Electronics for Future-Ready Smart Societies.” It provides a multidisciplinary forum for presenting original research, exchanging innovative ideas, and discussing advances across intelligent computing, next-generation communication systems, cyber-physical infrastructure, green electronics, robotics, healthcare technologies, and secure digital solutions. By encouraging interdisciplinary research and collaboration among academia, industry, and research institutions, INDICON 2026 seeks to promote technological innovation that addresses real-world challenges in sustainability, connectivity, security, and societal well-being. The conference is scheduled for December 18–20, 2026, at Sri Sairam Engineering College, Chennai, India. \cite{sarkar2023leaf}

The 23rd Edition of the IEEE India Council International Conference (INDICON 2026) brings together the global research community to explore emerging technologies and their role in shaping sustainable, intelligent, and future-ready societies. Organized by the IEEE India Council in association with the IEEE Madras Section, the conference is guided by the theme “Net-Zero Cyber-Physical Intelligence: AI, 6G \& Sustainable Electronics for Future-Ready Smart Societies.” It provides a multidisciplinary forum for presenting original research, exchanging innovative ideas, and discussing advances across intelligent computing, next-generation communication systems, cyber-physical infrastructure, green electronics, robotics, healthcare technologies, and secure digital solutions. By encouraging interdisciplinary research and collaboration among academia, industry, and research institutions, INDICON 2026 seeks to promote technological innovation that addresses real-world challenges in sustainability, connectivity, security, and societal well-being. The conference is scheduled for December 18–20, 2026, at Sri Sairam Engineering College, Chennai, India. The Table \ref{tab:tab1} is mentined here.

\begin{table}[htbp]
\caption{Table Type Styles}
\begin{center}
\begin{tabular}{|c|c|c|c|}
\hline
\textbf{Table}&\multicolumn{3}{|c|}{\textbf{Table Column Head}} \\
\cline{2-4} 
\textbf{Head} & \textbf{\textit{Table column subhead}}& \textbf{\textit{Subhead}}& \textbf{\textit{Subhead}} \\
\hline
copy& More table copy& &  \\
\hline
\end{tabular}
\label{tab:tab1}
\end{center}
\end{table}


The 23rd Edition of the IEEE India Council International Conference (INDICON 2026) brings together the global research community to explore emerging technologies and their role in shaping sustainable, intelligent, and future-ready societies. Organized by the IEEE India Council in association with the IEEE Madras Section, the conference is guided by the theme “Net-Zero Cyber-Physical Intelligence: AI, 6G \& Sustainable Electronics for Future-Ready Smart Societies.” It provides a multidisciplinary forum for presenting original research, exchanging innovative ideas, and discussing advances across intelligent computing, next-generation communication systems, cyber-physical infrastructure, green electronics, robotics, healthcare technologies, and secure digital solutions.  The sample code for reference Image \ref{fig:IEEEINDIA} is given. By encouraging interdisciplinary research and collaboration among academia, industry, and research institutions, INDICON 2026 seeks to promote technological innovation that addresses real-world challenges in sustainability, connectivity, security, and societal well-being. The conference is scheduled for December 18–20, 2026, at Sri Sairam Engineering College, Chennai, India.




\begin{figure}
    \centering
    \includegraphics[width=0.5\linewidth]{1.png}
    \caption{Example of a figure caption.}
    \label{fig:IEEEINDIA}
\end{figure}

The 23rd Edition of the IEEE India Council International Conference (INDICON 2026) brings together the global research community to explore emerging technologies and their role in shaping sustainable, intelligent, and future-ready societies. Please number citations consecutively within brackets \cite{zheng2006comprehensive}. Organized by the IEEE India Council in association with the IEEE Madras Section, the conference is guided by the theme “Net-Zero Cyber-Physical Intelligence: AI, 6G \& Sustainable Electronics for Future-Ready Smart Societies.” It provides a multidisciplinary forum for presenting original research, exchanging innovative ideas, and discussing advances across intelligent computing, next-generation communication systems, cyber-physical infrastructure, green electronics, robotics, healthcare technologies, and secure digital solutions. By encouraging interdisciplinary research and collaboration among academia, industry, and research institutions, INDICON 2026 seeks to promote technological innovation that addresses real-world challenges in sustainability, connectivity, security, and societal well-being. The conference is scheduled for December 18–20, 2026, at Sri Sairam Engineering College, Chennai, India.






\bibliographystyle{IEEEtran} 
\bibliography{References}

\end{document}
